\section{结果与讨论}

\subsection{模型性能评价}
本研究构建的模型在独立测试集上表现优异，能够较好地区分不同实验条件下纳米硒颗粒的毒性差异。
具体性能评估结合后续特征重要性分析、数据集统计描述与局限性讨论进行综合解释。

% \begin{figure}[ht]
%   \centering
%   \includegraphics[width=0.8\textwidth]{images/roc_curve.png}
%   \caption{模型预测的 ROC 曲线}
%   \label{fig:roc}
% \end{figure}

\subsection{特征重要性分析}
增强特征模型的特征重要性分析表明，IC50、元素 Se 含量、长期稳定性和结合能 Se3d 是影响硒纳米颗粒细胞毒性预测结果的核心变量。其中，IC50 直接反映剂量--效应关系，元素 Se 含量与颗粒纯度及潜在离子释放背景相关，长期稳定性表征颗粒在分散体系中的聚集与氧化趋势，结合能 Se3d 则反映表面化学态差异。该结果说明，在单一材料体系的毒性建模中，能够刻画材料本征化学属性的增强特征，比单纯依赖粒径、浓度和暴露时间等通用参数更具解释力。

\subsection{数据集特征统计描述}

\subsubsection{样本概况}
本研究数据集共包含 383 条有效记录（剔除粒径、Zeta 电位等关键理化性质及细胞存活率缺失记录后），
来源于 47 篇已发表文献，时间跨度为 2008--2019 年。其中 2018 年
贡献样本量最多（157 条，占 41.0\%），其次为 2019 年（96 条，
占 25.1\%）和 2012 年（44 条，占 11.5\%）。所有纳米硒颗粒的形貌
均为球形，物种来源以人类细胞为主（371 条，占 96.9\%）。

\subsubsection{连续型特征分布}
表 \ref{tbl:statistics} 总结了主要连续型特征的描述性统计。粒径（直径）的
中位数为 30 nm（四分位距 30--92 nm），均值（65.30 nm）显著高于中位数，
偏度为 3.92，呈明显右偏分布，表明存在少量大粒径样本（最大值 500 nm）。
水合粒径的中位数为 91 nm（四分位距 84--110 nm），普遍大于原始直径，
反映了颗粒在水相中的水化层效应。Zeta 电位的中位数为 $-7.50$ mV，范围
覆盖 $-60$ 至 $36.2$ mV，表明不同表面修饰对颗粒表面电荷的调控范围较广。
浓度分布右偏程度最高（偏度 8.28），中位数仅 0.90 $\mu$g/mL，均值为
12.59 $\mu$g/mL，最大值达 490 $\mu$g/mL，提示少数高浓度实验数据可能
对模型训练产生较大影响。暴露时间呈轻微左偏分布（偏度 $-0.42$），中位数
72 h，四分位距 24--72 h。细胞存活率分布较为均匀（偏度 $-0.17$），均值
49.43\%，中位数 50.00\%，正负样本比例相对均衡。

\begin{table}[htbp]
  \centering
  \caption{数据集连续型特征的描述性统计（各特征可用样本量见 $n$ 列）}
  \label{tbl:statistics}
  \begin{tabular}{lccccccc}
    \hline
    特征 & $n$ & 均值 & 标准差 & 中位数 & Q1 & Q3 & 偏度 \\
    \hline
    粒径 (nm)            & 339 & 65.30  & 79.72  & 30.00  & 30.00  & 92.15   & 3.92  \\
    水合粒径 (nm)        & 279 & 104.97 & 46.73  & 91.00  & 84.00  & 110.00  & 1.69  \\
    Zeta 电位 (mV)       & 144 & $-5.19$ & 29.76  & $-7.50$ & $-29.10$ & 29.10 & $-0.08$ \\
    暴露时间 (h)         & 363 & 55.21  & 26.14  & 72.00  & 24.00  & 72.00   & $-0.42$ \\
    浓度 ($\mu$g/mL)     & 363 & 12.59  & 38.17  & 0.90   & 0.52   & 10.00   & 8.28  \\
    细胞存活率 (\%)      & 371 & 49.43  & 28.50  & 50.00  & 41.00  & 68.00   & $-0.17$ \\
    \hline
  \end{tabular}
\end{table}

\subsubsection{类别型特征分布}
数据集涉及 42 种细胞系，以结肠癌细胞 HCT 116 出现频率最高
（98 次，25.6\%），其次为肝癌细胞 HepG2（41 次，10.7\%）、
黑色素瘤细胞 A375（23 次，6.0\%）和乳腺癌细胞 MCF-7（23 次，
6.0\%），前 10 种细胞系合计覆盖 71.0\% 的样本。组织来源以结肠
（28.7\%）、肝脏（16.7\%）和脑组织（8.9\%）为主。

表面修饰类型共 43 种，其中聚糖修饰（PTR）占比最高（127 次，
33.2\%），其次为聚糖蛋白复合物 PR（18 次，4.7\%）和
PAMAM+RGD+siRNA 复合修饰（15 次，3.9\%）。合成方法以化学
还原法为主（352 次，占 91.9\%），毒性检测方法以 MTT 法最为
常见（206 次，53.8\%）。

\subsection{数据局限性}
本研究存在以下局限性：（1）数据集规模有限（383 条有效记录），
来源于文献手工提取，约 40.6\% 的原始记录因粒径、Zeta 电位等关键理化性质及细胞存活率字段缺失而
被剔除，且可能存在发表偏倚，即阴性结果或无显著毒性的数据较少
被报道；（2）不同文献的实验条件（如毒性检测方法涉及 MTT、MTS、
CCK-8 等不同试剂）存在异质性，可能引入系统性噪声；
（3）数据集时间跨度为 2008--2019 年，早期文献的实验方法和报告
标准与近期存在差异。未来研究可通过多中心数据共享、标准化实验
方案以及主动学习策略来扩充数据规模并降低异质性。
